problem:  
Let \((M,g)\) be a complete, noncompact Riemannian manifold of dimension \(n\ge2\) that exhibits **two distinct diffusion regimes**:  
- (**Local Gaussian regime**) There exist constants \(C,c>0\) such that for \(t\le1\) and all \(x,y\in M\),
  \[
  C^{-1}t^{-n/2}\exp\!\left(-\frac{d(x,y)^2}{ct}\right)
  \le p_t(x,y)
  \le Ct^{-n/2}\exp\!\left(-\frac{d(x,y)^2}{C t}\right),
  \]
  i.e., standard parabolic behavior and local spectral dimension \(d_s^{\text{loc}}=n\).  

- (**Large-scale sub-Gaussian regime**) There exist parameters \(d_f^\infty>0, d_w^\infty>2\) and constants \(C',c'>0\) such that for \(t\ge1\) and all \(x,y\in M\),
  \[
  C'^{-1}t^{-d_f^\infty/d_w^\infty}
  \exp\!\left[-c'\!\left(\frac{d(x,y)^{d_w^\infty}}{t}\right)^{\!\frac{1}{d_w^\infty-1}}\right]
  \le p_t(x,y)
  \le C't^{-d_f^\infty/d_w^\infty}
  \exp\!\left[-\frac1{C'}\!\left(\frac{d(x,y)^{d_w^\infty}}{t}\right)^{\!\frac{1}{d_w^\infty-1}}\right].
  \]
  The corresponding large-scale spectral dimension is \(d_s^{\infty} = 2d_f^\infty / d_w^\infty\).

This two-regime structure can arise, for example, on manifolds that are approximately Euclidean near a compact core but possess ends with anomalous (e.g. fractal-like or degenerate) geometry that slows diffusion at infinity.

Consider the nonlinear heat equation
\[
u_t = \Delta_g u + u^p, \qquad u(\cdot,0) = u_0 \ge 0\in C_c^\infty(M),
\]
and let \(T(u_0)\) denote the maximal existence time.

---

**Problem (Two-scale Fujita crossover and escape-to-infinity blow-up):**  
1. Does there exist a critical exponent \(p_c(M,g)\) separating global existence and finite-time blow-up such that  
   \[
   1+\frac{2}{d_s^{\infty}} \;\le p_c(M,g)\;\le 1+\frac{2}{d_s^{\text{loc}}}\,?
   \]
   (Since the large-scale diffusion is slower, typically \(d_s^{\infty}<d_s^{\text{loc}}\), implying \(1+2/d_s^{\infty}>1+2/d_s^{\text{loc}}\).)

2. Can there be an intermediate range \(p\in(1+2/d_s^{\text{loc}},\,1+2/d_s^{\infty})\) in which *finite-time blow-up occurs but migrates to infinity*, i.e. for which \(T(u_0)<\infty\) yet
   \[
   \sup_{K\Subset M} |u(x,t)|<\infty \text{ for all } t<T(u_0),
   \quad
   \text{and}\quad \limsup_{t\uparrow T(u_0)}\sup_{d(x,x_0)>\!R}|u(x,t)|\to\infty
   \text{ for every } R>0?
   \]
   Such blow-up would be driven by the geometry-induced slowdown of diffusion at infinity.

3. Is there a functional or geometric invariant—e.g., involving harmonic measure of infinity, renormalized spectral dimension, or return probability tail—that exactly determines \(p_c(M,g)\) in mixed-scale settings?

---

**Heuristic insight:**  Local diffusion tends to suppress blow-up (strong smoothing), while large-scale anomalous diffusion at infinity weakens dissipation. A crossover exponent should thus interpolate between the local and asymptotic Fujita thresholds; the transition could produce new singularity types where energy escapes toward infinity faster than it accumulates locally.

---

Why is it a "good" problem:  
This question probes a **new frontier in nonlinear geometric analysis**: how nonlinear blow-up behavior depends on the **competition between distinct diffusion scales** on a single manifold. Classical Fujita theory relies on a single diffusion dimension, but here one must analyze how small-time Gaussian heat propagation interacts with large-time sub-Gaussian slowdown. The problem unifies:
- **Geometric analysis and heat kernel theory**, through rigorous two-regime estimates;
- **Nonlinear PDE blow-up theory**, via scaling and concentration mechanisms;
- **Probabilistic geometry**, since diffusion transition corresponds to different random-walk regimes.

Its apparently simple question—does geometry with two diffusion regimes yield a *crossover Fujita exponent* or *escape-to-infinity blow-up*—conceals deep analytic structure. A solution would introduce new invariants measuring how geometry at infinity alters nonlinear thresholds, creating a bridge between global geometric diffusion and nonlinear singularity formation.  This synthesis, novelty, and potential impact make it a genuine “good’’ research problem.